This computational example compares three forms of free motion of a damped
mass--spring system. All three cases start with the same displacement and zero
velocity. Only the damping coefficient changes.

\section{Model and initial conditions}
The displacement $x(t)$, measured from equilibrium, satisfies
\begin{equation}
 m\ddot{x}+c\dot{x}+kx=0,
 \qquad x(0)=1\,\mathrm{m},\quad \dot{x}(0)=0.
\end{equation}
We use $m=1\,\mathrm{kg}$ and $k=4\,\mathrm{N/m}$, so that
\begin{equation}
 \omega_0=\sqrt{k/m}=2\,\mathrm{rad/s},
 \qquad \zeta=\frac{c}{2\sqrt{mk}}.
\end{equation}
The damping coefficient $c$ is measured in $\mathrm{N\,s/m}$.

\begin{center}
\begin{tabular}{lrr}
Case & $c\;(\mathrm{N\,s/m})$ & $\zeta$ \\
\hline
Underdamped & $0.8$ & $0.2$ \\
Critically damped & $4.0$ & $1.0$ \\
Overdamped & $8.0$ & $2.0$ \\
\end{tabular}
\end{center}

\section{Published results}
\begin{center}
\includegraphics[width=0.95\linewidth]{comparison.png}
\end{center}
The underdamped mass repeatedly crosses equilibrium. For these initial
conditions, critical damping approaches equilibrium without overshoot and
faster than the overdamped case. Increasing damping beyond the critical value
slows the return.

The mechanical energy is
\begin{equation}
 E(t)=\frac{1}{2}m\dot{x}^{2}+\frac{1}{2}kx^{2},
 \qquad \frac{dE}{dt}=-c\dot{x}^{2}\leq 0.
\end{equation}
Displacement can oscillate while mechanical energy decreases. The curves are
sampled at $601$ times from $0$ to $12\,\mathrm{s}$, with spacing
$0.02\,\mathrm{s}$; this spacing is an output sampling interval, not an
integration step size.

\section{Reproducing the calculation in Julia}
Writing $\mathbf{u}=(x,\dot{x})^{T}$ gives
\begin{equation}
 \mathbf{u}(t)=\exp(tA)\mathbf{u}(0),\qquad
 A=\begin{pmatrix}0&1\\-k/m&-c/m\end{pmatrix}.
\end{equation}
The supplied Julia script evaluates the matrix exponential with the standard
\texttt{LinearAlgebra} library, using double-precision arithmetic. Tests compare
the displacement and velocity with independent scalar analytic solutions in
all three damping regimes, and check energy decay.

The \PMlinkexternal{preset viewer and downloadable Julia project}
{https://images.physicslibrary.org/examples/julia-oscillator/index.html}
include the source, environment manifest, CSV data, plots, and provenance.
Extract the project archive and run these commands in its top-level directory:
\begin{verbatim}
julia --startup-file=no --project=. test/runtests.jl
julia --startup-file=no --project=. oscillator.jl
\end{verbatim}
The calculation uses Julia standard libraries. To explore different initial
conditions locally, change \texttt{presets.toml} and rerun the tests and script.
The published tests describe the original presets; adjust the preset-specific
expectations if changing the physical experiment.

\section{Interpretation and limits}
This model assumes a linear spring and viscous damping, with no driving force.
It does not model dry friction, nonlinear springs, or large-deformation effects.
The displayed results are saved computations. Viewing the article or changing
the displayed preset does not start a calculation on the PhysicsLibrary server.
